\chapter{Σχετικές Εργασίες}
\label{ch:sota}


\section{Από την εξαγωγή χαρακτηριστικών στα \tl{CNNs} και τους
\tl{Vision Transformers}}
\label{sec:related-cnn-vt}

Στα πρώτα σημαντικά βήματα για την εξέλιξη της υπολογιστικής όρασης
συγκαταλέγονται προσεγγίσεις που βασίζονταν στον χειροκίνητο σχεδιασμό
χαρακτηριστικών από εικόνες (\tl{handcrafted features}) για την περιγραφή του οπτικού
περιεχομένου. Τα χαρακτηριστικά αυτά εξάγονταν από τις εικόνες και στη
συνέχεια αξιοποιούνταν από αλγορίθμους αναγνώρισης. Σημαντικές μέθοδοι
προς αυτή την κατεύθυνση αποτέλεσαν το
\tl{Scale-Invariant Feature Transform (SIFT)}
\cite{lowe2004distinctive}, το οποίο εξάγει τοπικά χαρακτηριστικά
ανθεκτικά σε μεταβολές της κλίμακας και της περιστροφής, και το
\tl{Histogram of Oriented Gradients (HOG)}
\cite{dalal2005histograms}, το οποίο περιγράφει τη μορφή και τη δομή των
αντικειμένων μέσω της κατανομής των κατευθύνσεων των κλίσεων
(\tl{gradients}) και εφαρμόστηκε αρχικά στην ανίχνευση ανθρώπων σε εικόνες.

Αντίστοιχες προσεγγίσεις αναπτύχθηκαν στη συνέχεια και για την αναγνώριση
ενεργειών σε βίντεο, με χαρακτηριστικό παράδειγμα να αποτελεί η μέθοδος
\tl{Improved Dense Trajectories} \cite{wang2013action}, η οποία αξιοποιεί
την κίνηση σημείων μεταξύ διαδοχικών καρέ και περιορίζει την επίδραση της
κίνησης της κάμερας, ώστε να αποτυπώνεται αποτελεσματικότερα η ανθρώπινη κίνηση.

Σημείο καμπής στην εξέλιξη της υπολογιστικής όρασης αποτέλεσε η αξιοποίηση
των νευρωνικών δικτύων, καθώς η εξαγωγή χαρακτηριστικών άρχισε πλέον να
μαθαίνεται αυτόματα από τα δεδομένα, περιορίζοντας την ανάγκη χειροκίνητου
σχεδιασμού. Σημαντικό ρόλο σε αυτή τη μετάβαση διαδραμάτισαν τα βαθιά
συνελικτικά νευρωνικά δίκτυα
(\tl{Convolutional Neural Networks -- CNNs}). Στα δίκτυα αυτά, τα βάρη κάθε συνελικτικού επιπέδου οργανώνονται σε
μικρά εκπαιδεύσιμα φίλτρα, τα οποία ονομάζονται πυρήνες συνέλιξης
(\tl{kernels}). Οι πυρήνες εφαρμόζονται σε διαφορετικές περιοχές της
εισόδου του επιπέδου μέσω της πράξης της συνέλιξης και προσαρμόζονται
κατά την εκπαίδευση, ώστε να ανιχνεύουν χρήσιμα χαρακτηριστικά. Με τη
διαδοχική εφαρμογή πολλών συνελικτικών επιπέδων, τα πρώτα επίπεδα
μαθαίνουν απλά τοπικά χαρακτηριστικά, όπως ακμές και υφές, ενώ τα
βαθύτερα τα συνδυάζουν σχηματίζοντας πιο σύνθετες αναπαραστάσεις
\cite{goodfellow2016deep}. Σημαντικό ορόσημο στην εξέλιξη των
\tl{CNNs} και στην ταξινόμηση εικόνων αποτέλεσε το \tl{AlexNet}
\cite{krizhevsky2012imagenet}.

Η δυνατότητα αυτόματης εκμάθησης χαρακτηριστικών επεκτάθηκε αργότερα
στην κατανόηση βίντεο και ειδικότερα στην αναγνώριση ανθρώπινων
ενεργειών, μέσω αρχιτεκτονικών όπως τα \tl{Two-Stream Networks}
\cite{simonyan2014two}, το \tl{C3D} \cite{tran2015learning} και το
\tl{I3D} \cite{carreira2017quo}. Τα \tl{Two-Stream Networks} συνδυάζουν
δύο \tl{CNNs}, ένα που επεξεργάζεται μεμονωμένα καρέ και ένα που
επεξεργάζεται την οπτική ροή (\tl{optical flow}) μεταξύ διαδοχικών
καρέ, ενώ τα \tl{C3D} και \tl{I3D} αξιοποιούν τρισδιάστατες
(\tl{3D}) συνελίξεις, οι οποίες εκτείνονται τόσο στον χώρο όσο και στον
χρόνο. Συνολικά, η αξιοποίηση των \tl{CNNs} καθόρισε σε μεγάλο βαθμό
την πορεία της έρευνας και συνέβαλε στη σημαντική βελτίωση των επιδόσεων
τόσο στην ταξινόμηση εικόνων όσο και στην αναγνώριση ανθρώπινων
ενεργειών.

Καθοριστικό ρόλο στην πρόοδο αυτή διαδραμάτισε η διαθεσιμότητα μεγάλων
επισημασμένων συνόλων δεδομένων. Στην αναγνώριση εικόνων, το
\tl{ImageNet} \cite{deng2009imagenet}, με περισσότερες από ένα
εκατομμύριο επισημασμένες εικόνες, κατέστησε εφικτή την εκπαίδευση
μοντέλων όπως το \tl{AlexNet}. Στην αναγνώριση ανθρώπινων ενεργειών, τα
\tl{HMDB51} \cite{kuehne2011hmdb} και \tl{UCF101}
\cite{soomro2012ucf101}, τα οποία χρησιμοποιούνται και στην παρούσα
εργασία, αποτέλεσαν για χρόνια βασικά σύνολα
αναφοράς, παρά το σχετικά περιορισμένο μέγεθός τους. Στα μεγαλύτερα
σύνολα δεδομένων περιλαμβάνεται το \tl{Kinetics-400}
\cite{kay2017kinetics}, το οποίο διαθέτει περίπου 300.000 βίντεο από 400
κλάσεις και χρησιμοποιείται ευρέως για την εκπαίδευση και την αξιολόγηση
μοντέλων για την αναγνώριση ανθρώπινων ενεργειών. Αντίστοιχα, το
\tl{Something-Something V2} \cite{goyal2017something} περιλαμβάνει
περίπου 220.000 βίντεο από 174 κλάσεις και χρησιμοποιείται κυρίως για την
αξιολόγηση της κατανόησης της χρονικής εξέλιξης ενεργειών και
αλληλεπιδράσεων με αντικείμενα. Λόγω των υψηλών υπολογιστικών απαιτήσεων που συνεπάγεται ο μεγάλος όγκος
τους, τα σύνολα αυτά δεν αξιοποιήθηκαν στην παρούσα εργασία, όπως
εξηγείται αναλυτικότερα στην Ενότητα~\ref{sec:dataset-usage}.

Τέλος, οι \tl{Vision Transformers} αποτελούν σήμερα μία από τις πλέον
σύγχρονες και ευρέως αξιοποιούμενες αρχιτεκτονικές στην υπολογιστική
όραση. Βασίζονται στον \tl{Transformer} \cite{vaswani2017attention} και
στον μηχανισμό της αυτοπροσοχής (\tl{self-attention}). Σε αντίθεση με
τα \tl{CNNs}, όπου τα χαρακτηριστικά προκύπτουν από την τοπική
συσχέτιση γειτονικών περιοχών της εισόδου, η αυτοπροσοχή επιτρέπει σε
κάθε τμήμα της εισόδου να αλληλεπιδρά άμεσα με όλα τα υπόλοιπα.
Προτάθηκαν αρχικά για την ταξινόμηση εικόνων
\cite{dosovitskiy2021image} και στη συνέχεια επεκτάθηκαν στην
επεξεργασία βίντεο \cite{arnab2021vivit,bertasius2021space}.

\section[Αυτοεπιβλεπόμενη προεκπαίδευση μοντέλων βίντεο]
{Αυτοεπιβλεπόμενη μάθηση ως μέθοδος\\
προεκπαίδευσης μοντέλων κατανόησης βίντεο}
\label{sec:related-self-supervised-learning}

Για πολλά χρόνια, η κύρια τεχνική εκπαίδευσης των μοντέλων υπολογιστικής
όρασης, όπως τα \tl{CNNs} που αναφέρθηκαν, ήταν η επιβλεπόμενη μάθηση
(\tl{supervised learning}). Βασικός της περιορισμός είναι ότι, για να
επιτευχθεί ικανοποιητική απόδοση, απαιτείται μεγάλος όγκος επισημασμένων
δεδομένων, η δημιουργία των οποίων αποτελεί χρονοβόρα και δαπανηρή
διαδικασία. Η αυτοεπιβλεπόμενη μάθηση (\tl{self-supervised learning}) ως
μέθοδος προεκπαίδευσης μοντέλων κατανόησης βίντεο έχει γνωρίσει
σημαντική ανάπτυξη τα τελευταία χρόνια, καθώς περιορίζει σημαντικά αυτή
την ανάγκη. Ωστόσο, πριν από την εφαρμογή της σε μοντέλα εικόνας
και βίντεο, η βασική της λογική είχε ήδη αξιοποιηθεί στην επεξεργασία
φυσικής γλώσσας (\tl{natural language processing -- NLP}). 
Στην προσέγγιση \tl{word2vec} \cite{mikolov2013efficient}, οι αναπαραστάσεις
των λέξεων μαθαίνονται μέσω της πρόβλεψης μιας λέξης από τις γειτονικές
της. Στη συνέχεια, το γλωσσικό μοντέλο
\tl{BERT} \cite{devlin2018bert} εισήγαγε την πρόβλεψη τυχαία
αποκρυμμένων λέξεων μιας πρότασης (\tl{masked language modeling}) και
αποτέλεσε ορόσημο για την προεκπαίδευση γλωσσικών μοντέλων, καθώς και μία
από τις πιο επιδραστικές εφαρμογές της αυτοεπιβλεπόμενης μάθησης. Η ιδέα
της απόκρυψης και πρόβλεψης τμημάτων της εισόδου μεταφέρθηκε αργότερα
και στην υπολογιστική όραση.

Αρχικά, για την προεκπαίδευση μοντέλων εικόνας χρησιμοποιήθηκαν βοηθητικές
εργασίες (\tl{pretext tasks}), όπως η ανακατασκευή της σωστής διάταξης
ανακατεμένων τμημάτων μιας εικόνας (\tl{jigsaw puzzle})
\cite{noroozi2016unsupervised} και η πρόβλεψη της γωνίας περιστροφής της
\cite{gidaris2018unsupervised}. Αντίστοιχα, στα δεδομένα βίντεο, το
μοντέλο εκπαιδεύτηκε να αναγνωρίζει αν μια ακολουθία καρέ βρίσκεται στη
σωστή χρονική σειρά \cite{misra2016shuffle}. Στη συνέχεια, αναπτύχθηκαν μέθοδοι αντιπαραθετικής μάθησης
(\tl{contrastive learning}), όπως οι \tl{MoCo} \cite{he2020momentum} και
\tl{SimCLR} \cite{chen2020simple}, οι οποίες
εκπαιδεύουν το μοντέλο να παράγει παρόμοιες αναπαραστάσεις για
διαφορετικές επαυξημένες εκδοχές της ίδιας εικόνας και να τις απομακρύνει 
από τις αναπαραστάσεις διαφορετικών εικόνων
στον χώρο των αναπαραστάσεων.
Ακολούθησαν μη αντιπαραθετικές προσεγγίσεις, όπως οι \tl{BYOL}
\cite{grill2020bootstrap} και \tl{DINO} \cite{caron2021emerging}, στις οποίες το μοντέλο εκπαιδεύεται
ώστε οι αναπαραστάσεις που παράγει να προσεγγίζουν εκείνες ενός
μοντέλου-στόχου, χωρίς να απαιτείται η χρήση αρνητικών δειγμάτων.

Στην πορεία, σημαντικό μέρος της έρευνας στράφηκε στην ανακατασκευή
αποκρυμμένων περιοχών της εισόδου στον χώρο των \tl{pixels}, μεταφέροντας
στην υπολογιστική όραση την ιδέα του \tl{BERT}. Στο \tl{Masked Autoencoder (MAE)}
\cite{he2022masked}, μεγάλο ποσοστό των \tl{patches} μιας εικόνας
αποκρύπτεται και το μοντέλο εκπαιδεύεται να ανακατασκευάσει τα \tl{pixels}
τους. Η προσέγγιση αυτή επεκτάθηκε στα δεδομένα βίντεο με το
\tl{VideoMAE} \cite{tong2022videomae} και στη συνέχεια με το
\tl{VideoMAE V2} \cite{wang2023videomaev2}, στο οποίο προτείνεται μια
στρατηγική διπλής απόκρυψης (\tl{dual masking}) για την αποδοτικότερη
προεκπαίδευση μοντέλων μεγάλης κλίμακας. Ειδικότερα, εφαρμόζεται απόκρυψη
τόσο στον κωδικοποιητή όσο και στον αποκωδικοποιητή. Ο κωδικοποιητής
επεξεργάζεται μόνο ένα μικρό ποσοστό των χωροχρονικών \tl{tokens} του
βίντεο, ενώ ο αποκωδικοποιητής ανακατασκευάζει στον χώρο των \tl{pixels}
ένα επιλεγμένο υποσύνολο των αποκρυμμένων περιοχών, με στόχο τη μείωση των
απαιτήσεων σε υπολογιστική ισχύ και μνήμη.

Σε συνέχεια του \tl{MAE}, αναπτύχθηκε μια διαφορετική κατεύθυνση, στην
οποία η έμφαση μετατοπίζεται από την ανακατασκευή του οπτικού περιεχομένου
στην πρόβλεψη απευθείας στον χώρο των αναπαραστάσεων. Στην κατεύθυνση αυτή
εντάσσονται οι αρχιτεκτονικές \tl{JEPA}. Η προσέγγιση αυτή υλοποιήθηκε για εικόνες με το \tl{I-JEPA}
\cite{mahmoud2023i-jepa} και επεκτάθηκε στη συνέχεια στο βίντεο με το
\tl{V-JEPA} \cite{adrien2023v-jepa}. Τα αποτελέσματα του \tl{V-JEPA}
δείχνουν ότι η πρόβλεψη χαρακτηριστικών οδηγεί σε ισχυρότερες
αναπαραστάσεις από τις μεθόδους ανακατασκευής \tl{pixels}, όταν οι παράμετροι του \tl{backbone} παραμένουν παγωμένες κατά την
εκπαίδευση του ταξινομητή. Η
δυνατότητα άμεσης αξιοποίησης των αναπαραστάσεων σε διαφορετικές εργασίες,
χωρίς περαιτέρω προσαρμογή του \tl{backbone}, αναδεικνύει την ικανότητα
γενίκευσής τους.

Η κλιμάκωση της αυτοεπιβλεπόμενης προεκπαίδευσης σε ολοένα μεγαλύτερες
και πιο ποικιλόμορφες συλλογές βίντεο συνέβαλε στην ανάπτυξη των μοντέλων
θεμελίωσης βίντεο (\tl{video foundation models})
\cite{bommasani2021opportunities}. Τα μοντέλα αυτά λειτουργούν ως γενικοί
κωδικοποιητές βίντεο (\tl{video encoders}), καθώς οι αναπαραστάσεις που μαθαίνουν
μπορούν να αξιοποιηθούν σε ένα ευρύ φάσμα εργασιών (\tl{downstream tasks}).

Εξέλιξη της παραπάνω προσέγγισης αποτελεί το \tl{V-JEPA 2}, το οποίο
κλιμακώνει την αυτοεπιβλεπόμενη προεκπαίδευση σε μεγαλύτερο όγκο
δεδομένων εικόνας και βίντεο, με στόχο όχι μόνο την κατανόηση αλλά και την
πρόβλεψη της εξέλιξης γεγονότων στον φυσικό κόσμο. Επιπλέον, η επέκταση
\tl{V-JEPA 2-AC}\cite{assran2025vjepa2} αξιοποιείται για τον σχεδιασμό ενεργειών ρομποτικών
βραχιόνων, όπως ο χειρισμός αντικειμένων (\tl{robotic planning}). Η
εργασία αυτή αναδεικνύει τη σταδιακή μετάβαση των μοντέλων θεμελίωσης
βίντεο από την κατανόηση οπτικού περιεχομένου προς τη μοντελοποίηση και
την αλληλεπίδραση με τον φυσικό κόσμο.


\section{Τεχνικές συμπίεσης μοντέλων}
\label{sec:related-efficient-vision-models}

Η αύξηση του υπολογιστικού κόστους, που συνοδεύει τη συνεχή κλιμάκωση
αυτών των μοντέλων, έχει καταστήσει αναγκαία την ανάπτυξη τεχνικών
συμπίεσης για την αποδοτικότερη εκτέλεσή τους \cite{li2023modelcompression}. Ενδεικτικά, το \tl{AlexNet} περιλαμβάνει περίπου 60 εκατομμύρια
παραμέτρους \cite{krizhevsky2012imagenet}, ενώ τα \tl{ViT-L/16} και
\tl{ViT-H/16 backbones} του \tl{V-JEPA} περιλαμβάνουν περίπου 300 και
630 εκατομμύρια παραμέτρους, αντίστοιχα
\cite{adrien2023v-jepa}. Παράλληλα, οι μεγαλύτερες εκδοχές των \tl{VideoMAE V2} και
\tl{V-JEPA 2} φτάνουν την κλίμακα του ενός δισεκατομμυρίου παραμέτρων
\cite{wang2023videomaev2,assran2025vjepa2}. Στη σχετική βιβλιογραφία, 
στις βασικές τεχνικές συμπίεσης περιλαμβάνονται: η αφαίρεση πλεοναζόντων παραμέτρων (\tl{pruning}), η κβαντοποίηση
(\tl{quantization}) και η απόσταξη γνώσης
(\tl{knowledge distillation})\cite{li2023modelcompression}. 

\subsection{\tl{Pruning}}
\label{subsec:pruning}

Η ιδέα της αφαίρεσης παραμέτρων των οποίων η απουσία επηρεάζει ελάχιστα
την απόδοση ενός δικτύου διατυπώθηκε ήδη από τις αρχές της δεκαετίας του
1990 \cite{lecun1990optimal}. Στην εργασία των \tl{Han et al.}
\cite{han2015learning}, αφαιρούνται τα βάρη με τη μικρότερη απόλυτη τιμή
και το δίκτυο επανεκπαιδεύεται, μειώνοντας τον αριθμό των παραμέτρων του
\tl{AlexNet} περίπου κατά εννέα φορές, χωρίς απώλεια ακρίβειας.

Η αφαίρεση μεμονωμένων βαρών, ωστόσο, δεν οδηγεί πάντοτε σε πραγματική
επιτάχυνση. Για τον λόγο αυτό αναπτύχθηκαν δομημένες μέθοδοι
(\tl{structured pruning}), στις οποίες αφαιρούνται ολόκληρα δομικά
στοιχεία του δικτύου, όπως φίλτρα ή κανάλια \cite{li2017pruning}. Στους
\tl{Vision Transformers} έχουν προταθεί επίσης μέθοδοι οι οποίες, αντί
για παραμέτρους, αφαιρούν προοδευτικά τα λιγότερο πληροφοριακά
\tl{tokens} κατά την επεξεργασία τους, μειώνοντας με αυτόν τον τρόπο το
υπολογιστικό κόστος \cite{rao2021dynamicvit}.

\subsection{\tl{Model quantization}}
\label{subsec:related-quantization}

Παράλληλα με το \tl{pruning}, η κβαντοποίηση
(\tl{model quantization}) αναπτύχθηκε ως μια ακόμη τεχνική μείωσης των
απαιτήσεων μνήμης και του υπολογιστικού κόστους των νευρωνικών δικτύων.
Κατά την εφαρμογή της, τα βάρη ή/και οι ενεργοποιήσεις ενός μοντέλου
μετατρέπονται από αριθμούς κινητής υποδιαστολής 32 \tl{bits} σε
αναπαραστάσεις χαμηλότερης ακρίβειας, όπως ακέραιες αναπαραστάσεις
8 \tl{bits} \cite{jacob2018quantization}, ενώ νεότερες μέθοδοι
περιορίζουν περαιτέρω την αριθμητική ακρίβεια, φτάνοντας ακόμη και στα
4 \tl{bits} \cite{li2023repqvit}. Η κβαντοποίηση μπορεί είτε να λαμβάνεται υπόψη κατά την εκπαίδευση,
ώστε το μοντέλο να προσαρμόζεται στη μειωμένη αριθμητική ακρίβεια
(\tl{quantization-aware training}), είτε να εφαρμόζεται μετά την
ολοκλήρωση της εκπαίδευσης (\tl{post-training quantization})
\cite{gholami2021survey}.

Η πρακτική αποδοτικότητα της κβαντοποίησης υπόκειται, ωστόσο, σε
περιορισμούς που σχετίζονται με το \tl{hardware} στο οποίο εκτελείται
το μοντέλο. Για να μειωθεί πραγματικά ο χρόνος εκτέλεσης, το
\tl{hardware} πρέπει να υποστηρίζει πράξεις στη μειωμένη αριθμητική
ακρίβεια. Όταν τέτοια υποστήριξη δεν είναι διαθέσιμη, μπορεί να
χρησιμοποιηθεί προσομοιωμένη κβαντοποίηση
(\tl{simulated quantization}), κατά την οποία οι τιμές περιορίζονται
στις διακριτές στάθμες της επιλεγμένης ακρίβειας, αλλά οι υπολογισμοί
εξακολουθούν να εκτελούνται σε αριθμούς κινητής υποδιαστολής. Με αυτόν
τον τρόπο αξιολογείται η επίδραση της κβαντοποίησης στην ακρίβεια του
μοντέλου, χωρίς όμως να επιτυγχάνεται πραγματική επιτάχυνση. Επιπλέον,
η υποστήριξη πράξεων χαμηλής ακρίβειας διαφέρει μεταξύ συσκευών,
ιδιαίτερα για ακρίβειες χαμηλότερες των 8 \tl{bits}, με αποτέλεσμα το
όφελος της ίδιας μεθόδου ως προς την ταχύτητα εκτέλεσης και την
ενεργειακή κατανάλωση να διαφέρει από συσκευή σε συσκευή
\cite{gholami2021survey}.

Στην κατηγορία της κβαντοποίησης μετά την εκπαίδευση ανήκει το
\tl{RepQ-ViT} \cite{li2023repqvit}, μια μέθοδος σχεδιασμένη για
αρχιτεκτονικές \tl{Vision Transformer}. Η μέθοδος αντιμετωπίζει τις
ιδιαίτερες κατανομές των ενεργοποιήσεων μετά τα επίπεδα \tl{LayerNorm}
και \tl{Softmax}, εφαρμόζοντας αρχικά κβαντοποίηση ανά κανάλι
(\tl{channel-wise quantization}) και λογαριθμική κβαντοποίηση,
αντίστοιχα. Στη συνέχεια, μέσω επαναπαραμετροποίησης των κλιμάκων
κβαντοποίησης (\tl{scale reparameterization}), οι μορφές αυτές 
μετατρέπονται σε ισοδύναμες, καταλληλότερες για αποδοτική εκτέλεση στο \tl{hardware}
χωρίς να απαιτείται εκ νέου εκπαίδευση ολόκληρου του μοντέλου.



\subsection{Προσεγγίσεις \tl{Knowledge Distillation}}
\label{subsec:related-knowledge-distillation}

Η γενική ιδέα της μεταφοράς γνώσης από ένα νευρωνικό δίκτυο σε ένα άλλο
διατυπώθηκε για πρώτη φορά το 1991 από τον \tl{Jürgen Schmidhuber}
\cite{schmidhuber1992learning}. Η διαδικασία περιγραφόταν ως σύμπτυξη της ιεραρχίας σε ένα ενιαίο
αναδρομικό δίκτυο (\tl{collapsing the hierarchy into a single recurrent
net}). Στο πλαίσιο της συμπίεσης μοντέλων, η ιδέα εφαρμόστηκε πιο άμεσα το
2006 από τους \tl{Bucilu\u{a} et al.} \cite{bucilua2006model}, οι
οποίοι συμπίεσαν ένα σύνολο μοντέλων (\tl{ensemble}) σε ένα μοναδικό,
μικρότερο δίκτυο. Για την εκπαίδευσή του χρησιμοποιήθηκαν μη
επισημασμένα πραγματικά δεδομένα ή συνθετικά δείγματα, στα οποία οι
προβλέψεις του συνόλου χρησιμοποιήθηκαν ως στόχοι.

Στη συνέχεια, οι \tl{Ba} και \tl{Caruana} \cite{ba2014deep} έδειξαν ότι
ρηχά δίκτυα μπορούν να προσεγγίσουν την απόδοση βαθύτερων μοντέλων,
μιμούμενα τις τελικές εξόδους τους (\tl{logits}). Η απόσταξη γνώσης πήρε τη σύγχρονη μορφή της με την εργασία των
\tl{Hinton et al.} \cite{hinton2015distilling}, η οποία καθιέρωσε τον
συγκεκριμένο όρο και τη χρήση της παραμέτρου θερμοκρασίας στη συνάρτηση
\tl{softmax}. Για τιμές μεγαλύτερες της μονάδας, ο συντελεστής αυτός
εξομαλύνει την κατανομή πιθανοτήτων του \tl{teacher} και αναδεικνύει
περισσότερη πληροφορία για τις σχέσεις μεταξύ των κλάσεων. Την ίδια περίοδο, το \tl{FitNets}
\cite{romero2015fitnets} επέκτεινε τη μεταφορά γνώσης πέρα από τις
τελικές εξόδους του \tl{teacher}, αξιοποιώντας και τις ενδιάμεσες
αναπαραστάσεις του.

Η απόσταξη γνώσης εφαρμόστηκε στη συνέχεια και στις αρχιτεκτονικές
\tl{Vision Transformers}. Στην εργασία \tl{Training Data-Efficient Image Transformers \&
Distillation Through Attention} \cite{touvron2021deit} παρουσιάζεται η
προσέγγιση \tl{DeiT} για την αποδοτική εκπαίδευση μοντέλων
\tl{Vision Transformer}. Στην εκδοχή της
που αξιοποιεί την απόσταξη γνώσης, εισάγεται στην είσοδο του \tl{student} ένα επιπλέον εκπαιδεύσιμο \tl{token}
(\tl{distillation token}). Τα επίπεδα
του \tl{Transformer} επεξεργάζονται το \tl{token} αυτό μαζί με το
\tl{class token}. Στην έξοδο του μοντέλου, η αναπαράσταση του
\tl{class token} χρησιμοποιείται για την πρόβλεψη των πραγματικών
ετικετών, ενώ η αναπαράσταση του \tl{distillation token} καθοδηγείται από
τις προβλέψεις του \tl{teacher}. Με αυτόν τον τρόπο, η επίβλεψη από τις
ετικέτες και η γνώση του \tl{teacher} συνδυάζονται κατά την εκπαίδευση του
\tl{student}.

Βασικό σημείο αναφοράς για την παρούσα εργασία αποτελεί η εργασία
\tl{Knowledge Distillation: A Good Teacher Is Patient and Consistent}
\cite{beyer2022knowledge}, στην οποία η απόσταξη γνώσης αντιμετωπίζεται ως
πρόβλημα αντιστοίχισης συναρτήσεων (\tl{function matching}). Οι συγγραφείς
δείχνουν ότι καθοριστικός παράγοντας για την αποτελεσματικότητα της
απόσταξης είναι η συνέπεια των εισόδων: ο \tl{teacher} και ο \tl{student}
πρέπει να δέχονται ακριβώς την ίδια επαυξημένη είσοδο. Στην προσέγγιση
\tl{FunMatch}, εφαρμόζεται επιπλέον η τεχνική επαύξησης \tl{mixup}, κατά
την οποία δύο εικόνες αναμειγνύονται σε μία, και ο \tl{student}
εκπαιδεύεται να προσεγγίζει την πρόβλεψη του \tl{teacher} για τη
συγκεκριμένη είσοδο. Παράλληλα, αναδεικνύεται η σημασία ενός προγράμματος εκπαίδευσης μεγάλης
διάρκειας για την αποτελεσματική μεταφορά της γνώσης.

Πέρα από τη μεταφορά των τελικών προβλέψεων, νεότερες μέθοδοι για
\tl{Vision Transformers} εξετάζουν την απόσταξη ενδιάμεσων
αναπαραστάσεων και σχέσεων μεταξύ των \tl{tokens}. Το \tl{ViTKD}
\cite{yang2024vitkd} εφαρμόζει άμεση αντιστοίχιση στα ρηχά επίπεδα,
ενώ στα βαθύτερα ανακατασκευάζει την αναπαράσταση του \tl{teacher} από
μερικώς αποκρυμμένα \tl{tokens} του \tl{student}. Αντίστοιχα, το \tl{TinyMIM} \cite{ren2023tinymim} δείχνει ότι
αποτελεσματικότερο στόχο απόσταξης αποτελούν οι σχέσεις που σχηματίζονται
μεταξύ των εσωτερικών αναπαραστάσεων των \tl{tokens} μέσω του μηχανισμού
προσοχής, αντί του \tl{class token} ή της άμεσης αντιστοίχισης
χαρακτηριστικών. Παράλληλα, το \tl{G2SD}
\cite{huang2023generic} συνδυάζει σε δύο διαδοχικά στάδια την απόσταξη γενικών
αναπαραστάσεων με τη μεταφορά των τελικών προβλέψεων για την
εργασία-στόχο.

Χαρακτηριστικό παράδειγμα εφαρμογής της απόσταξης σε μεγαλύτερη κλίμακα
αποτελεί το \tl{DINOv2} \cite{oquab2023dinov2}. Οι συγγραφείς
προεκπαιδεύουν αρχικά ένα μοντέλο \tl{ViT-g} μεγάλης κλίμακας και στη
συνέχεια το χρησιμοποιούν ως παγωμένο \tl{teacher} για την απόσταξη της
γνώσης του σε μικρότερες αρχιτεκτονικές \tl{ViT}. Στην περίπτωση αυτή, η απόσταξη
δεν περιορίζεται στις προβλέψεις μιας συγκεκριμένης εργασίας, αλλά
μεταφέρει τις γενικεύσιμες αναπαραστάσεις που αποκτήθηκαν κατά την
αυτοεπιβλεπόμενη προεκπαίδευση. Η λογική αυτή παρουσιάζει άμεση
συνάφεια με την απόσταξη του \tl{backbone} στην παρούσα εργασία.

Αντίστοιχες προσεγγίσεις έχουν εφαρμοστεί και στην προεκπαίδευση
μοντέλων βίντεο. Στη μέθοδο \tl{MVD} \cite{wang2023mvd}, αναπαραστάσεις
προεκπαιδευμένων μοντέλων εικόνας και βίντεο χρησιμοποιούνται ως στόχοι
για την πρόβλεψη αποκρυμμένων χωροχρονικών περιοχών, με τα πρώτα να μεταφέρουν ισχυρότερη χωρική και τα δεύτερα ισχυρότερη
χρονική πληροφορία.
Στο \tl{UMT} \cite{li2023unmasked}, η απόσταξη χρησιμοποιείται για την
εκπαίδευση ενός \tl{Video Transformer}, έχοντας ως \tl{teacher} ένα
προεκπαιδευμένο μοντέλο θεμελίωσης εικόνας. Ο \tl{student}
επεξεργάζεται τα ορατά \tl{tokens} ενός μερικώς αποκρυμμένου βίντεο και
ευθυγραμμίζει τις αναπαραστάσεις τους με τους αντίστοιχους
σημασιολογικούς στόχους που παράγει ο \tl{teacher}.

Συνοψίζοντας, η απόσταξη γνώσης έχει εξεταστεί μέσω της μεταφοράς των
τελικών προβλέψεων \cite{touvron2021deit,beyer2022knowledge}, των
αναπαραστάσεων και των σχέσεων μεταξύ των \tl{tokens}
\cite{yang2024vitkd,ren2023tinymim,oquab2023dinov2}, καθώς και μέσω του
συνδυασμού διαφορετικών στόχων \cite{huang2023generic}. Ειδικότερα, οι
προσεγγίσεις \tl{DeiT}, \tl{FunMatch}, \tl{ViTKD}, \tl{TinyMIM},
\tl{G2SD} και \tl{DINOv2} αναπτύχθηκαν και αξιολογήθηκαν κυρίως σε
μοντέλα και εργασίες του πεδίου των εικόνων. Αντίθετα, στα δεδομένα
βίντεο, οι μέθοδοι \tl{MVD} και \tl{UMT} αξιοποιούν την απόσταξη κυρίως
για την προεκπαίδευση νέων μοντέλων και όχι για τη συμπίεση ενός ήδη
προεκπαιδευμένου μοντέλου θεμελίωσης βίντεο. Η παρούσα εργασία εξετάζει αυτή την κατεύθυνση, συγκρίνοντας τη μεταφορά
τελικών προβλέψεων, αναπαραστάσεων και τον συνδυασμό τους από το
\tl{V-JEPA Large} στις μικρότερες παραλλαγές \tl{Base}, \tl{Small} και
\tl{Tiny}.  Οι παραλλαγές αυτές διαφέρουν ως προς το μέγεθος του \tl{encoder} που
χρησιμοποιείται ως \tl{backbone} και διαθέτουν σημαντικά λιγότερες
παραμέτρους από το \tl{V-JEPA Large}. Στόχος είναι να διερευνηθεί σε
ποιο βαθμό μπορούν, μέσω της απόσταξης γνώσης, να διατηρήσουν την
απόδοση του \tl{teacher} με χαμηλότερες υπολογιστικές απαιτήσεις.